\chapter{同态、正规子群与商群}
\label{chap:homomorphism-quotient}

两个群的元素名称和外观可以完全不同，却可能具有相同的复合规律。同态只保留乘法结构，因此可以把一个复杂群压缩到较简单的像；被压缩成单位元的那部分构成核。

若想把核中的元素视为彼此不可区分，必须保证等价类的乘法不依赖代表元。这个要求正是正规性。商群不是把符号 $G/H$ 写出来就自动存在，而是乘法良定义后才得到的新群；第一同构定理随后说明，按核取商与直接看像其实是同一结构。

\section{共轭、类与正规子群}
\label{sec:conjugacy-normal}

群论里一个看似奇怪、但物理意义非常直接的组合是
\[
xgx^{-1}.
\]
如果 $g$ 是某个操作，而 $x$ 表示“改变观察坐标或重新标记对象”，那么先用 $x^{-1}$ 把问题搬到旧标记下，做 $g$，再用 $x$ 搬回来，得到的正是同一个几何或代数结构在新标记下的表达。因此 $g$ 与 $xgx^{-1}$ 应被视为同一种操作的不同位置或不同坐标表达。这个关系叫共轭。

例如在正方形对称群中，绕一条对角线反射和绕另一条对角线反射可以通过整体转动 $90^\circ$ 相互变换，所以它们属于同一个共轭类。后面特征标为什么在一整个共轭类上取同一个值，本质上就是因为表示矩阵在共轭下只发生相似变换，而迹不随相似变换改变。


\begin{definition}[共轭]
若 $x,y\in G$，并且存在 $g\in G$ 使
\[
y=gxg^{-1},
\]
则称 $x$ 与 $y$ 共轭，记为 $x\sim y$。
\label{def:conjugacy}
\end{definition}

\begin{proposition}[共轭是等价关系]
群上的共轭关系满足自反性、对称性和传递性，因此把 $G$ 分割成若干互不相交的共轭类。
\label{prop:conjugacy-equivalence}

\end{proposition}

\begin{proof}
自反性来自 $x=exe^{-1}$。若 $y=gxg^{-1}$，则
\[
x=g^{-1}yg,
\]
故有对称性。若 $y=gxg^{-1}$ 且 $z=hyh^{-1}$，则
\[
z=(hg)x(hg)^{-1},
\]
故有传递性。
\end{proof}

\begin{definition}[共轭类、中心与中心化子]
元素 $x\in G$ 的共轭类为
\[
\operatorname{Cl}(x)=\{gxg^{-1}\mid g\in G\}.
\]
$G$ 的中心定义为
\[
Z(G)=\{z\in G\mid zg=gz,\ \forall g\in G\}.
\]
$x$ 的中心化子定义为
\[
C_G(x)=\{g\in G\mid gx=xg\}.
\]
\label{def:center-centralizer}
\end{definition}

\begin{proposition}[中心化子是子群]
对任意 $x\in G$，$C_G(x)\leqslant G$；并且
\[
x\in Z(G)
\quad\Longleftrightarrow\quad
C_G(x)=G
\quad\Longleftrightarrow\quad
\operatorname{Cl}(x)=\{x\}.
\]
\label{prop:centralizer-subgroup}

\end{proposition}

\begin{proof}
若 $a,b\in C_G(x)$，则 $ax=xa$、$bx=xb$，从而
\[
(ab^{-1})x=a(b^{-1}x)=a(xb^{-1})=(ax)b^{-1}=x(ab^{-1}).
\]
故 $ab^{-1}\in C_G(x)$，由一步子群判据得到 $C_G(x)\leqslant G$。其余等价关系直接由定义得到。
\end{proof}

\begin{theorem}[共轭类大小公式]
若 $G$ 为有限群，则
\[
\lvert\operatorname{Cl}(x)\rvert=[G:C_G(x)]
=\frac{\lvert G\rvert}{\lvert C_G(x)\rvert}.
\label{eq:class-size}
\]
特别地，每个共轭类的大小都整除 $\lvert G\rvert$。
\label{thm:class-size}

\end{theorem}

\begin{proof}
定义从左陪集集合到共轭类的映射
\[
\Phi:\{gC_G(x)\mid g\in G\}\longrightarrow \operatorname{Cl}(x),
\qquad
\Phi(gC_G(x))=gxg^{-1}.
\]
中心化子 $C_G(x)$ 一般并不正规，因此此处需要的只是它的\emph{左陪集集合}。即使很多文献仍把这个集合简写为 $G/C_G(x)$，也不应误以为它自动带有商群乘法。
先证良定义。若 $gC_G(x)=hC_G(x)$，则 $h^{-1}g\in C_G(x)$，所以
\[
(h^{-1}g)x=x(h^{-1}g).
\]
左右适当乘以 $h$ 与 $g^{-1}$ 得
\[
gxg^{-1}=hxh^{-1}.
\]
因此同一陪集给出同一共轭元。任意 $y\in\operatorname{Cl}(x)$ 都按定义可写为 $y=gxg^{-1}$，故映射满射。若
\[
gxg^{-1}=hxh^{-1},
\]
则 $h^{-1}g\in C_G(x)$，故 $gC_G(x)=hC_G(x)$，所以 $\Phi$ 单射。于是共轭类与左陪集一一对应，结合 Lagrange 定理即得\eqnrefs{eq:class-size}。
\end{proof}

这个公式其实就是轨道--稳定子定理的一个特例。让 $G$ 作用在自己上，作用定义为
\[
x\cdot g=xgx^{-1}.
\]
那么 $g$ 的轨道正是共轭类 $\operatorname{Cl}(g)$；使 $g$ 不动的那些 $x$ 满足 $xgx^{-1}=g$，等价于 $xg=gx$，所以稳定子正是中心化子 $C_G(g)$。于是
\[
|\operatorname{Cl}(g)|=\frac{|G|}{|C_G(g)|}.
\]

这给出一个非常有用的物理直觉：一个操作与“越多”的群元交换，它的中心化子越大，它的共轭类就越小；反之，一个操作越特殊、越少有别的操作与它兼容，它在群作用下会产生越大的等价轨道。

在分子对称性中，这条公式的直接读法是：共轭类大小恰好等于“与某类对称操作等价的操作的个数”。以 $D_3$（即 $C_{3v}$ 分子，如 $\mathrm{NH_3}$）为例，三个反射 $\sigma_v$ 在群作用下彼此等价（旋转把一条镜面转到另一条），共轭类大小为 $3$；而两个旋转 $C_3,C_3^2$ 也彼此等价（反射把逆时针转成顺时针），共轭类大小为 $2$。后续用特征标投影算符构造对称适配线性组合（对称性适配线性组合）时，每个共轭类贡献一个独立的投影通道，因此共轭类的个数等于不可约表示的个数——这正是为什么本章必须把共轭类结构算清楚：它直接决定了下一章特征标表的行数和列数。

\begin{corollary}[类方程]
设 $G$ 为有限群，从每个非平凡共轭类中取一个代表元 $x_1,\dots,x_r$，则
\[
\lvert G\rvert
=
\lvert Z(G)\rvert
+
\sum_{i=1}^{r}[G:C_G(x_i)].
\label{eq:class-equation}
\]
\label{cor:class-equation}

\end{corollary}

\begin{proof}
共轭关系把 $G$ 分割成互不相交的共轭类。中心中的每个元素 $z\in Z(G)$ 都满足 $gzg^{-1}=z$，所以其共轭类只有一个元素；反过来，若某元素的共轭类只有一个元素，则它与所有 $g\in G$ 交换，因此属于 $Z(G)$。于是所有单元素共轭类恰好组成 $Z(G)$。其余共轭类取代表 $x_i$，由类大小定理
\[
|\operatorname{Cl}(x_i)|=[G:C_G(x_i)].
\]
把所有互不相交共轭类的大小相加，便得到
\[
|G|=|Z(G)|+\sum_{i=1}^r[G:C_G(x_i)].
\]
\end{proof}

类方程把整个群拆成"中心"与"非中心共轭类"两部分，是分析有限群结构最基本的计数工具。它的一个直接推论是：若 $|G|$ 是素数幂 $p^n$（$n\ge1$），则 $Z(G)$ 非平凡。理由是每个非中心共轭类的大小 $[G:C_G(x_i)]$ 既整除 $|G|=p^n$ 又大于 $1$，故都是 $p$ 的倍数；于是类方程右端除 $|Z(G)|$ 外各项都被 $p$ 整除，而左端 $p^n$ 也被 $p$ 整除，因此 $|Z(G)|$ 必须被 $p$ 整除，特别 $|Z(G)|\ge p>1$。这一结论是 $p$ 群结构定理的起点，也是后续证明有限 $p$ 群必有非平凡正规子群的基石。对 $D_3$ 而言，类方程 $6=1+2+3$ 中 $|Z(D_3)|=1$ 不被任何素数整除，这与 $6$ 不是素数幂一致。

\medskip
\noindent\emph{$D_3$ 共轭类的几何解读。}
在 $D_3$ 中，共轭的几何含义是"换一个观察角度后，同一个操作长什么样"。用生成元 $r,s$ 可以把三个共轭类完全读出来，而不必依赖置换的循环型。

旋转类 $\{r,r^2\}$。
关系 $srs^{-1}=r^2$（即\eqnrefs{eq:d3-srs-relation}）直接说明 $r$ 与 $r^2$ 共轭。几何上，先反射、再转 $120^\circ$、再反射回去，等于转 $240^\circ$：反射把逆时针转角翻成了顺时针转角。因此 $r$ 与 $r^2$ 是同一种操作（"转 $120^\circ$"）在镜像观察下的两种写法，必须归入同一共轭类。中心化子 $C_{D_3}(r)=\left\langle r\right\rangle$（只有旋转与旋转交换），故由\cref{thm:class-size}
\[
\lvert\operatorname{Cl}(r)\rvert=\frac{6}{3}=2,
\]
恰好容纳 $r$ 与 $r^2$。

反射类 $\{s,sr,sr^2\}$。
三个反射分别对应三条反射轴（过顶点 $1$、过顶点 $2$、过顶点 $3$ 的中线）。旋转 $r$ 把一条轴转到下一条，因此用 $r$ 做共轭就在三个反射之间循环：
\[
rsr^{-1}=r s r^2=(r s)r^2=(sr^2)r^2=sr^4=sr.
\]
于是 $s$ 与 $sr$ 共轭；再共轭一次得 $sr$ 与 $sr^2$ 共轭，三个反射同属一类。中心化子 $C_{D_3}(s)=\left\langle s\right\rangle$（只有 $e$ 和 $s$ 本身与 $s$ 交换），故
\[
\lvert\operatorname{Cl}(s)\rvert=\frac{6}{2}=3,
\]
恰好容纳三个反射。

单位元类 $\{e\}$。
$e$ 与所有群元交换，故 $geg^{-1}=e$，自成一类。

于是 $D_3$ 的三个共轭类为 $\{e\}$、$\{r,r^2\}$、$\{s,sr,sr^2\}$，大小 $1,2,3$，中心 $Z(D_3)=\{e\}$。类方程验证：
\[
6=1+2+3.
\]
这个例子揭示了一条贯穿全书的物理直觉：\emph{同一共轭类中的操作在物理上不可区分}——换一个对称操作作为观察基准，一个旋转就变成它的逆，一条反射轴就变成另一条。后面特征标之所以在整类上取同一个值，正是因为特征标度量的是"操作的物理痕迹"，而物理痕迹不依赖观察角度。
\par\medskip

陪集已经把群切成若干块，于是一个很自然的想法是：能不能把"每一个陪集"当成一个新的群元，并定义
\[
(gH)(kH)=(gk)H?
\]
如果这件事可行，我们就把原群中属于同一个陪集的细节全部压缩掉，得到一个更小的群。这就是商群思想。

难点在于，一个陪集可以用很多不同代表元来写。例如 $gH=g'H$、$kH=k'H$ 时，我们希望无论选 $g,k$ 还是 $g',k'$，乘积陪集都一样；否则上式不是一个真正的二元运算。恰恰在这里，普通子群不够，必须要求 $H$ 在所有共轭下保持自己，也就是正规性。换句话说，“正规子群”不是为了多记一种子群，而是为了让陪集乘法良定义。

\begin{definition}[正规子群]
若 $N\leqslant G$ 且对任意 $g\in G$ 都有
\[
gNg^{-1}=N,
\]
则称 $N$ 为 $G$ 的正规子群，记作
\[
N\mathrel{\trianglelefteq} G.
\]
\label{def:normal-subgroup}
\end{definition}

“正规”并不是额外的闭合条件，而是说子群在整个群的共轭作用下保持不变。因此正规子群恰好是能够被压缩成“商群中的一个单位元”的子群。

\begin{theorem}[正规性的等价刻画]
设 $N\leqslant G$。下列条件等价：
\begin{enumerate}[label=(\roman*),leftmargin=3em]
  \item $N\mathrel{\trianglelefteq} G$；
  \item 对任意 $g\in G$，$gN=Ng$；
  \item 对任意 $g\in G$，$gNg^{-1}\subseteq N$。
\end{enumerate}
\label{thm:normal-equivalences}

\end{theorem}

\begin{proof}
$(\mathrm{i})\Rightarrow(\mathrm{ii})$：由 $gNg^{-1}=N$ 右乘 $g$ 得 $gN=Ng$。

$(\mathrm{ii})\Rightarrow(\mathrm{iii})$：由 $gN=Ng$ 右乘 $g^{-1}$ 得 $gNg^{-1}=N$，因而当然包含于 $N$。

$(\mathrm{iii})\Rightarrow(\mathrm{i})$：把条件用于 $g^{-1}$ 得
\[
g^{-1}Ng\subseteq N.
\]
左右分别乘以 $g$ 与 $g^{-1}$，得到
\[
N\subseteq gNg^{-1}.
\]
与原包含关系合并即为相等。
\end{proof}

\begin{proposition}[两个基本正规子群]
对任意群 $G$，有
\[
\{e\}\mathrel{\trianglelefteq} G,
\qquad
G\mathrel{\trianglelefteq} G,
\qquad
Z(G)\mathrel{\trianglelefteq} G.
\]
若 $G$ 为 Abel 群，则每个子群都正规。
\label{prop:basic-normal-subgroups}

\end{proposition}

\begin{proof}
对子群 $\{e\}$，任意 $g\in G$ 都有
\[
g\{e\}g^{-1}=\{geg^{-1}\}=\{e\},
\]
故其正规。对 $G$ 本身，共轭映射 $x\mapsto gxg^{-1}$ 是 $G$ 到自身的双射，所以 $gGg^{-1}=G$。若 $z\in Z(G)$，则 $gzg^{-1}=z\in Z(G)$，因此
\[
gZ(G)g^{-1}\subseteq Z(G).
\]
把同一论证用于 $g^{-1}$ 得反向包含，故 $Z(G)\mathrel{\trianglelefteq} G$。

若 $G$ 为 Abel 群，任取子群 $H\leqslant G$ 与 $g\in G$，对每个 $h\in H$ 有
\[
ghg^{-1}=hgg^{-1}=h,
\]
所以 $gHg^{-1}=H$。因此每个子群都正规。
\end{proof}

\begin{definition}[商群]
设 $N\mathrel{\trianglelefteq} G$。所有陪集组成的集合
\[
G/N=\{gN\mid g\in G\}
\]
在乘法
\[
(gN)(hN)=(gh)N
\label{eq:quotient-product}
\]
下构成群，称为 $G$ 对 $N$ 的商群。
\label{def:quotient-group}
\end{definition}

上式最需要证明的并不是结合律，而是\textbf{良定义性}：同一陪集可以选不同代表元，结果必须与代表元的选择无关。

\begin{theorem}[商群乘法的良定义性]
若 $N\mathrel{\trianglelefteq} G$，且
\[
gN=g'N,
\qquad
hN=h'N,
\]
则
\[
(gh)N=(g'h')N.
\]
因此\eqnrefs{eq:quotient-product} 不依赖代表元的选择。
\label{thm:quotient-well-defined}

\end{theorem}

\begin{proof}
由陪集相等判据，存在 $n_1,n_2\in N$ 使
\[
g'=gn_1,
\qquad
h'=hn_2.
\]
于是
\begin{align*}
g'h'
&=gn_1hn_2\\
&=gh\,(h^{-1}n_1h)n_2.
\end{align*}
因 $N\mathrel{\trianglelefteq} G$，有 $h^{-1}n_1h\in N$，从而
\[
(h^{-1}n_1h)n_2\in N.
\]
故 $g'h'\in ghN$，即 $(g'h')N=(gh)N$。
\end{proof}

这里“良定义”三个字必须真正理解。一个陪集不是一个单独元素，而是一整组代表元。例如 $gN=gnN$ 对任意 $n\in N$ 都成立。如果我们写
\[
(gN)(hN)=ghN,
\]
就必须保证把 $g$ 换成同一陪集里的 $gn_1$、把 $h$ 换成 $hn_2$ 后，结果仍是同一个陪集。计算
\[
(gn_1)(hn_2)=gh\,(h^{-1}n_1h)n_2.
\]
若 $N\mathrel{\trianglelefteq} G$，则 $h^{-1}n_1h\in N$，于是括号中的两个因子都在 $N$，它们的乘积也在 $N$，最终仍得到 $ghN$。正规性正是在这一步被使用的。

如果 $N$ 不是正规子群，这一步可能失败：不同代表元会给出不同“乘积陪集”，那么所谓商群乘法根本不是函数。由此可以看到，正规子群的定义并不是形式上的附加条件，而是商结构能够存在的必要机制。


\begin{corollary}[有限商群的阶]
若 $G$ 有限且 $N\mathrel{\trianglelefteq} G$，则
\[
\lvert G/N\rvert=[G:N]=\frac{\lvert G\rvert}{\lvert N\rvert}.
\]
\label{cor:quotient-order}

\end{corollary}

\begin{proof}
商群 $G/N$ 的元素就是 $N$ 在 $G$ 中的左陪集，因此其元素个数按定义等于指数 $[G:N]$。Lagrange 定理给出
\[
|G|=[G:N]|N|,
\]
两边除以 $|N|$ 即得
\[
|G/N|=[G:N]=\frac{|G|}{|N|}.
\]
\end{proof}

\begin{example}[$D_3$ 的正规子群与商群：旋转子群为何特殊]
\label{ex:d3-normal-subgroup-geometry}
$D_3$ 的旋转子群 $\left\langle r\right\rangle=\{e,r,r^2\}$ 是正规的，而三个反射子群 $\left\langle s\right\rangle,\left\langle sr\right\rangle,\left\langle sr^2\right\rangle$ 都不正规。这两件事的几何原因都可以从共轭直接读出。

$\left\langle r\right\rangle$ 正规。反射共轭旋转给出其逆：$srs^{-1}=r^2\in\left\langle r\right\rangle$，$sr^2s^{-1}=r\in\left\langle r\right\rangle$。几何上，"先反射、转 $120^\circ$、再反射回去"等于反向转 $120^\circ$，仍是旋转。旋转共轭旋转当然仍是旋转。因此 $\left\langle r\right\rangle$ 在所有共轭下保持自身，$g\left\langle r\right\rangle g^{-1}=\left\langle r\right\rangle$。

$\left\langle s\right\rangle$ 不正规。旋转共轭反射给出另一个反射：$rsr^{-1}=sr\neq s$，而 $sr\notin\left\langle s\right\rangle=\{e,s\}$。几何上，旋转把过顶点 $1$ 的反射轴转到过顶点 $2$ 的反射轴，因此共轭把 $\left\langle s\right\rangle$ 里的反射 $s$ 送到属于另一个子群 $\left\langle sr\right\rangle$ 的反射 $sr$。于是 $r\left\langle s\right\rangle r^{-1}=\left\langle sr\right\rangle\neq\left\langle s\right\rangle$。

因为 $\left\langle r\right\rangle\mathrel{\trianglelefteq} D_3$，可以构造商群。两个陪集为
\[
\left\langle r\right\rangle=\{e,r,r^2\},
\qquad
s\left\langle r\right\rangle=\{s,sr,sr^2\},
\]
分别代表"纯旋转"与"含反射"。商群乘法
\[
(s\left\langle r\right\rangle)^2=s^2\left\langle r\right\rangle=\left\langle r\right\rangle
\]
说明两个反射合成一个旋转，因此
\[
D_3/\left\langle r\right\rangle\cong C_2.
\]
物理上，这个商群把 $C_{3v}$ 的六操作压缩成"真转动"与"含镜面反射"两类，恰好对应行列式 $\det R=\pm1$ 的区分。在分子光谱中，这种"偶/奇"分类决定了某些跃迁选择定则的符号：真转动保持手性，含反射的操作翻转手性。
\end{example}


正规子群与商群也最好通过一个具体计算消化。$S_3$ 中的偶置换组成
\[
A_3=\{e,(123),(132)\}.
\]
它是指数为 $2$ 的子群，因此自动正规。商群只有两个陪集：
\[
A_3,\qquad (12)A_3=\{(12),(13),(23)\}.
\]
把第二个陪集记作 $xA_3$，则
\[
(xA_3)^2=x^2A_3=A_3,
\]
因为任意换位 $x$ 都满足 $x^2=e$。所以 $S_3/A_3$ 只有单位元和一个二阶元素，因而
\[
S_3/A_3\cong C_2.
\]
物理上可以把这个商群理解为：如果我们完全忽略“具体是哪一个偶置换”的区别，只保留“偶/奇”两类信息，$S_3$ 就被压缩成了 $C_2$。

\section{同态、同构与自同构}
\label{sec:homomorphism-isomorphism}

前面一直在同一个群内部研究子群、陪集和共轭。现在要比较两个不同的群。比较群不能只比较“元素长什么样”，因为同一个抽象结构可能用完全不同的对象实现：正三角形的六个几何对称与三个对象的六个置换就是典型例子。真正应该保留的是乘法规律。

因此我们研究映射
\[
\varphi:G\longrightarrow H
\]
时，最基本的要求不是它把某个特定元素送到哪里，而是它必须满足
\[
\varphi(g_1g_2)=\varphi(g_1)\varphi(g_2).
\]
左边是“先在 $G$ 中相乘再映射”，右边是“先分别映射到 $H$ 再相乘”。二者一致意味着映射没有破坏群结构，这就是同态。若映射还一一且满，就说明两个群只是元素名字不同，结构完全相同，这就是同构。


\begin{definition}[群同态]
设 $G,H$ 为群。映射
\[
\varphi:G\longrightarrow H
\]
若满足
\[
\varphi(g_1g_2)=\varphi(g_1)\varphi(g_2),
\qquad
\forall g_1,g_2\in G,
\label{eq:homomorphism}
\]
则称 $\varphi$ 为群同态。
\label{def:homomorphism}
\end{definition}

同态的定义\textbf{不要求}单射或满射。满射同态称为满同态；双射同态才是同构。这样定义的好处是核、像与商群之间的结构关系可以在最一般的情形下统一表述。

\begin{proposition}[同态保持单位元、逆元与整数次幂]
若 $\varphi:G\to H$ 为同态，则
\[
\varphi(e_G)=e_H,
\qquad
\varphi(g^{-1})=\varphi(g)^{-1},
\qquad
\varphi(g^n)=\varphi(g)^n,
\quad n\in\ZZ.
\]
\label{prop:homomorphism-basic}

\end{proposition}

\begin{proof}
由
\[
\varphi(e_G)=\varphi(e_Ge_G)=\varphi(e_G)^2
\]
在 $H$ 中左消去 $\varphi(e_G)$，得 $\varphi(e_G)=e_H$。再由
\[
e_H=\varphi(e_G)=\varphi(gg^{-1})=\varphi(g)\varphi(g^{-1})
\]
得到 $\varphi(g^{-1})=\varphi(g)^{-1}$。幂公式随后由归纳法及逆元公式得到。
\end{proof}

\begin{definition}[核与像]
对群同态 $\varphi:G\to H$，定义
\[
\operatorname{Ker}\varphi
=\{g\in G\mid \varphi(g)=e_H\},
\]
\[
\operatorname{Im}\varphi
=\{\varphi(g)\mid g\in G\}\subseteq H.
\]
\label{def:kernel-image}
\end{definition}

\begin{theorem}[核与像的结构]
若 $\varphi:G\to H$ 为群同态，则
\[
\operatorname{Ker}\varphi\mathrel{\trianglelefteq} G,
\qquad
\operatorname{Im}\varphi\leqslant H.
\]
并且
\[
\varphi\text{ 为单射}
\quad\Longleftrightarrow\quad
\operatorname{Ker}\varphi=\{e_G\}.
\]
\label{thm:kernel-image}

\end{theorem}

\begin{proof}
若 $a,b\in\operatorname{Ker}\varphi$，则
\[
\varphi(ab^{-1})
=\varphi(a)\varphi(b)^{-1}
=e_H,
\]
故 $ab^{-1}\in\operatorname{Ker}\varphi$，所以核为子群。进一步，对 $g\in G$、$k\in\operatorname{Ker}\varphi$，
\[
\varphi(gkg^{-1})
=\varphi(g)e_H\varphi(g)^{-1}
=e_H,
\]
故 $gkg^{-1}\in\operatorname{Ker}\varphi$，核正规。

若 $x=\varphi(a)$、$y=\varphi(b)$ 属于像，则
\[
xy^{-1}
=\varphi(a)\varphi(b)^{-1}
=\varphi(ab^{-1})\in\operatorname{Im}\varphi,
\]
故像为子群。

若 $\operatorname{Ker}\varphi=\{e_G\}$ 且 $\varphi(a)=\varphi(b)$，则
\[
\varphi(b^{-1}a)=e_H,
\]
于是 $b^{-1}a=e_G$，即 $a=b$，故 $\varphi$ 单射。反之，若 $\varphi$ 单射且 $k\in\operatorname{Ker}\varphi$，则 $\varphi(k)=e_H=\varphi(e_G)$，从而 $k=e_G$，故核平凡。
\end{proof}

\begin{theorem}[第一同构定理]
对任意群同态 $\varphi:G\to H$，存在自然同构
\[
G/\operatorname{Ker}\varphi
\cong
\operatorname{Im}\varphi.
\label{eq:first-isomorphism}
\]
若 $\varphi$ 满射，则特别有
\[
G/\operatorname{Ker}\varphi\cong H.
\]
\label{thm:first-isomorphism}

\end{theorem}

\begin{proof}
令 $K=\operatorname{Ker}\varphi$。定义
\[
\overline{\varphi}:G/K\longrightarrow\operatorname{Im}\varphi,
\qquad
\overline{\varphi}(gK)=\varphi(g).
\]
若 $gK=g'K$，则 $g'^{-1}g\in K$，于是
\[
\varphi(g'^{-1}g)=e_H
\Longrightarrow
\varphi(g')^{-1}\varphi(g)=e_H
\Longrightarrow
\varphi(g)=\varphi(g'),
\]
故 $\overline{\varphi}$ 良定义。它保持乘法：
\[
\overline{\varphi}((gK)(hK))
=\varphi(gh)
=\varphi(g)\varphi(h).
\]
它按像的定义满射。若 $\overline{\varphi}(gK)=e_H$，则 $g\in K$，所以 $gK=K$，从而核平凡，故 $\overline{\varphi}$ 单射。于是得到\eqnrefs{eq:first-isomorphism}。
\end{proof}

第一同构定理可以用“压缩信息”的图像理解。同态 $\varphi:G\to H$ 会把若干不同的 $G$ 中元素映到同一个像。哪些元素被识别在一起？如果 $\varphi(g_1)=\varphi(g_2)$，则
\[
\varphi(g_2^{-1}g_1)=e_H,
\]
所以 $g_2^{-1}g_1\in\operatorname{Ker}\varphi$，也就是 $g_1$ 与 $g_2$ 位于同一个核陪集。换言之，同态丢失的信息恰好由核描述。

因此先把 $G$ 按核的陪集压缩成 $G/\operatorname{Ker}\varphi$，就已经把所有会被 $\varphi$ 识别的元素合并了；压缩后不再有额外信息损失，于是得到的群与实际像 $\operatorname{Im}\varphi$ 同构：
\[
G/\operatorname{Ker}\varphi\cong\operatorname{Im}\varphi.
\]
这条定理以后会反复解释“为什么某个群除掉一个冗余子群后变成另一个群”，例如 $SU(2)/\{\pm I\}\cong SO(3)$。

第一同构定理有两个同样重要的同伴——第二和第三同构定理，它们描述商群的"约化"和"嵌入"规则，在分析群结构时与第一同构定理配合使用。

\begin{theorem}[第二同构定理]
\label{thm:second-isomorphism}
设 $H\le G$，$N\mathrel{\trianglelefteq} G$。则 $HN\le G$，$H\cap N\mathrel{\trianglelefteq} H$，且
\begin{equation}
HN/N\cong H/(H\cap N).
\label{eq:second-isomorphism}
\end{equation}

\end{theorem}

\begin{proof}
对 $h_1n_1,h_2n_2\in HN$，因 $N\mathrel{\trianglelefteq} G$，有 $n_1h_2=h_2 n'$（$n'=h_2^{-1}n_1h_2\in N$），故
\[
(h_1n_1)(h_2n_2)=h_1h_2 n'n_2\in HN.
\]
逆元 $(hn)^{-1}=n^{-1}h^{-1}=h^{-1}(h n^{-1}h^{-1})\in HN$（因 $h n^{-1}h^{-1}\in N$）。故 $HN\le G$。

为比较两个商群，定义
\[
\varphi:H\to HN/N,\qquad\varphi(h)=hN.
\]
这是包含映射 $H\hookrightarrow HN$ 后取商的复合，故是同态。它满射：对 $hnN\in HN/N$，因 $n\in N$ 故 $hnN=hN=\varphi(h)$。

其核为
\[
\operatorname{Ker}\varphi=\{h\in H\mid hN=N\}=\{h\in H\mid h\in N\}=H\cap N.
\]
$H\cap N\mathrel{\trianglelefteq} H$ 因 $N\mathrel{\trianglelefteq} G$ 限制到 $H$ 上仍正规。

由第一同构定理 $H/\operatorname{Ker}\varphi\cong\operatorname{Im}\varphi$，即 $H/(H\cap N)\cong HN/N$。
\end{proof}

\begin{theorem}[第三同构定理]
\label{thm:third-isomorphism}
设 $N\mathrel{\trianglelefteq} G$，$K\mathrel{\trianglelefteq} G$，且 $N\le K$。则 $K/N\mathrel{\trianglelefteq} G/N$，且
\begin{equation}
(G/N)/(K/N)\cong G/K.
\label{eq:third-isomorphism}
\end{equation}

\end{theorem}

\begin{proof}
对 $gN\in G/N$，$kN\in K/N$，
\[
(gN)(kN)(gN)^{-1}=gkg^{-1}N\in K/N,
\]
因 $gkg^{-1}\in K$（$K\mathrel{\trianglelefteq} G$）。

再定义
\[
\varphi:G/N\to G/K,\qquad\varphi(gN)=gK.
\]
良定义性来自 $gN=g'N\Rightarrow g'^{-1}g\in N\le K\Rightarrow gK=g'K$。同态性由 $\varphi(gN\cdot hN)=ghK=(gK)(hK)$ 得到；任意 $gK$ 都是 $\varphi(gN)$，故映射满射。

这个满同态的核为
\[
\operatorname{Ker}\varphi=\{gN\mid gK=K\}=\{gN\mid g\in K\}=K/N.
\]

第一同构定理于是给出 $(G/N)/\operatorname{Ker}\varphi\cong\operatorname{Im}\varphi$，即 $(G/N)/(K/N)\cong G/K$。
\end{proof}

\begin{theorem}[对应定理（子群对应）]
\label{thm:correspondence}
设 $N\mathrel{\trianglelefteq} G$，$\pi:G\to G/N$ 为自然投影。则映射
\[
\Phi:\{H\mid N\le H\le G\}\longrightarrow\{H'\mid H'\le G/N\},
\qquad
\Phi(H)=H/N=\pi(H)
\]
是子群格之间的保序双射，且保持正规性和商群结构：
\begin{enumerate}[label=(\roman*),leftmargin=3em]
\item $H_1\le H_2\iff\Phi(H_1)\le\Phi(H_2)$；
\item $H\mathrel{\trianglelefteq} G\iff\Phi(H)\mathrel{\trianglelefteq} G/N$，此时 $(G/N)/\Phi(H)\cong G/H$；
\item $[G:H]=[G/N:\Phi(H)]$（有限指数）。
\end{enumerate}

\end{theorem}

\begin{proof}
定义候选逆映射
\[
\Psi:\{H'\le G/N\}\to\{H\mid N\le H\le G\},
\qquad
\Psi(H')=\pi^{-1}(H').
\]
需验证 $\pi^{-1}(H')$ 含 $N$：因 $N=\operatorname{Ker}\pi$，$e_{G/N}\in H'$，故 $N=\pi^{-1}(e_{G/N})\subseteq\pi^{-1}(H')$。

直接计算可得
\[
\Psi(\Phi(H))=\pi^{-1}(H/N)=\{g\in G\mid gN\in H/N\}=H,
\]
因 $gN=hN$（某 $h\in H$）等价于 $g\in hN\subseteq H$（因 $N\le H$）。
\[
\Phi(\Psi(H'))=\pi(\pi^{-1}(H'))=H',
\]
因 $\pi$ 满射。

商映射保持包含关系，所以 $H_1\le H_2\Rightarrow H_1/N\le H_2/N$；反过来由逆对应 $\Psi$ 得 $H_1/N\le H_2/N\Rightarrow H_1\le H_2$。

$H\mathrel{\trianglelefteq} G\Rightarrow H/N\mathrel{\trianglelefteq} G/N$：对 $gN\in G/N$，$(gN)(H/N)(gN)^{-1}=gHg^{-1}/N=H/N$。反过来用 $\Psi$ 拉回。

商群同构 $(G/N)/(H/N)\cong G/H$ 是第三同构定理。指数公式由 $[G:H]=|G/H|=|(G/N)/(H/N)|=[G/N:H/N]$ 得到。
\end{proof}

\medskip
\noindent\emph{对应定理的应用：$D_4$ 的子群格。}
$D_4=\langle r,s\mid r^4=s^2=e,\,srs=r^{-1}\rangle$，$|D_4|=8$。取 $N=\langle r^2\rangle=\{e,r^2\}\mathrel{\trianglelefteq} D_4$（$r^2$ 与所有元素交换），则 $D_4/N\cong V_4$（Klein 四群，因 $D_4/N$ 是 $4$ 阶 Abel 群且非循环）。

$V_4$ 有 $5$ 个子群：$\{e\}$、三个 $C_2$、$V_4$ 本身。由对应定理，$D_4$ 中含 $N$ 的子群恰好 $5$ 个：$N$ 本身、三个 $4$ 阶子群 $\langle r\rangle$、$\langle r^2,s\rangle$、$\langle r^2,rs\rangle$、以及 $D_4$。$D_4$ 的其余子群（不含 $N$ 的）有 $\{e\}$、$\langle s\rangle$、$\langle rs\rangle$、$\langle r^2s\rangle$、$\langle r^3s\rangle$，共 $4$ 个。故 $D_4$ 共有 $9$ 个子群。

对应定理还给出正规性信息：$V_4$ 的三个 $C_2$ 子群都正规（因 $V_4$ 是 Abel 群），故 $D_4$ 中对应的三个 $4$ 阶子群都正规。这与 $D_4$ 的几何一致：$\langle r\rangle$（纯转动）、$\langle r^2,s\rangle$（含水平反射）、$\langle r^2,rs\rangle$（含对角反射）都是正规子群。
\par\medskip


当同态 $\varphi:G\to H$ 为双射时，称 $\varphi$ 为\emph{同构}，记作 $G\cong H$。同构的两个群在代数上完全相同，只是元素名称不同。

\begin{definition}[同构]
若群同态 $\varphi:G\to H$ 为双射，则称 $\varphi$ 为同构，并称 $G$ 与 $H$ 同构，记作
\[
G\cong H.
\]
\label{def:isomorphism}
\end{definition}

同构表示两个群具有相同的群论结构，只是群元的“具体实现”不同。群的阶、交换性、群元阶的分布、子群格结构以及共轭类结构等均在同构下保持。

\begin{proposition}[共轭子群同构]
若 $K=gHg^{-1}$，则映射
\[
c_g:H\to K,
\qquad
c_g(h)=ghg^{-1}
\]
是同构。
\label{prop:conjugate-subgroups-isomorphic}

\end{proposition}

\begin{proof}
对 $h_1,h_2\in H$，有
\[
c_g(h_1h_2)=gh_1h_2g^{-1}
=(gh_1g^{-1})(gh_2g^{-1})
=c_g(h_1)c_g(h_2),
\]
所以 $c_g$ 是同态。它按 $K=gHg^{-1}$ 的定义满射。若 $c_g(h_1)=c_g(h_2)$，左右分别乘以 $g^{-1}$ 与 $g$ 得 $h_1=h_2$，所以单射。等价地，其逆映射就是
\[
c_{g^{-1}}:K\to H,
\qquad k\mapsto g^{-1}kg.
\]
故 $c_g$ 是同构。
\end{proof}

\begin{definition}[自同构与内自同构]
群 $G$ 到自身的同构称为自同构。所有自同构在映射复合下构成群
\[
\operatorname{Aut}(G).
\]
对固定 $g\in G$，映射
\[
c_g:G\to G,
\qquad
c_g(x)=gxg^{-1}
\]
称为由 $g$ 诱导的内自同构。所有内自同构构成 $\operatorname{Aut}(G)$ 的子群，记为 $\operatorname{Inn}(G)$。
\label{def:automorphism-inner}
\end{definition}

\begin{theorem}[内自同构群的结构]
映射
\[
\Psi:G\longrightarrow\operatorname{Inn}(G),
\qquad
\Psi(g)=c_g
\]
是满同态，且
\[
\operatorname{Ker}\Psi=Z(G).
\]
因此由第一同构定理，
\[
\operatorname{Inn}(G)\cong G/Z(G).
\label{eq:inner-aut-quotient}
\]
\label{thm:inner-aut-quotient}

\end{theorem}

\begin{proof}
由
\[
c_g\circ c_h(x)
=g(hxh^{-1})g^{-1}
=(gh)x(gh)^{-1}
=c_{gh}(x)
\]
知 $\Psi(gh)=\Psi(g)\Psi(h)$，所以 $\Psi$ 为同态；其像按定义就是 $\operatorname{Inn}(G)$，故满射。又
\[
c_g=\id(G)
\quad\Longleftrightarrow\quad
gxg^{-1}=x,\ \forall x\in G
\quad\Longleftrightarrow\quad
g\in Z(G).
\]
故核为 $Z(G)$，再用第一同构定理得\eqnrefs{eq:inner-aut-quotient}。
\end{proof}

第一同构定理也可以立刻用整数群看懂。定义
\[
\varphi:\ZZ\to C_n,\qquad \varphi(k)=a^k,
\]
其中 $C_n=\langle a\mid a^n=e\rangle$。这是同态，因为
\[
\varphi(k+\ell)=a^{k+\ell}=a^ka^\ell=\varphi(k)\varphi(\ell).
\]
哪些整数被映到单位元？恰好是 $n$ 的整数倍：
\[
\operatorname{Ker}\varphi=n\ZZ.
\]
而像就是整个 $C_n$。第一同构定理因此给出
\[
\ZZ/n\ZZ\cong C_n.
\]
这里左边的商群把所有相差 $n$ 的整数视为同一个剩余类；右边则把这些剩余类写成 $e,a,\dots,a^{n-1}$。这正是模 $n$ 算术与循环群之间的结构等价。


\cref{thm:inner-aut-quotient} 告诉我们 $\operatorname{Inn}(G)\cong G/Z(G)$，但 $\operatorname{Aut}(G)$ 通常比 $\operatorname{Inn}(G)$ 大。商群 $\operatorname{Aut}(G)/\operatorname{Inn}(G)$ 称为\textbf{外自同构群} $\operatorname{Out}(G)$，它衡量 $G$ 有多少"非内"的对称性。

\begin{proposition}[$\operatorname{Inn}(G)\mathrel{\trianglelefteq}\operatorname{Aut}(G)$]
内自同构群 $\operatorname{Inn}(G)$ 是 $\operatorname{Aut}(G)$ 的正规子群。
\end{proposition}

\begin{proof}
对任意 $\varphi\in\operatorname{Aut}(G)$ 和内自同构 $c_g\in\operatorname{Inn}(G)$，需证 $\varphi\circ c_g\circ\varphi^{-1}\in\operatorname{Inn}(G)$。对任意 $x\in G$：
\[
(\varphi\circ c_g\circ\varphi^{-1})(x)=\varphi(g\varphi^{-1}(x)g^{-1})=\varphi(g)\,x\,\varphi(g)^{-1}=c_{\varphi(g)}(x).
\]
故 $\varphi\circ c_g\circ\varphi^{-1}=c_{\varphi(g)}\in\operatorname{Inn}(G)$。
\end{proof}

\begin{theorem}[循环群的自同构群]
\label{thm:cyclic-aut}
设 $C_n=\langle a\rangle$ 是 $n$ 阶循环群。则
\[
\operatorname{Aut}(C_n)\cong(\ZZ/n\ZZ)^\times,
\]
其中 $(\ZZ/n\ZZ)^\times$ 是模 $n$ 的乘法可逆元群（Euler $\varphi$ 函数给出其阶 $|\operatorname{Aut}(C_n)|=\varphi(n)$）。具体地，每个自同构形如 $\sigma_k:a\mapsto a^k$，其中 $\gcd(k,n)=1$。
\end{theorem}

\begin{proof}
先自同构的形式。设 $\sigma\in\operatorname{Aut}(C_n)$。因 $a$ 生成 $C_n$，$\sigma$ 由 $\sigma(a)$ 完全确定。设 $\sigma(a)=a^k$（$0\le k<n$），则 $\sigma(a^m)=a^{km}$。

随后满射条件。$\sigma$ 是自同构当且仅当 $a^k$ 也生成 $C_n$，即 $\langle a^k\rangle=C_n$。这等价于 $\gcd(k,n)=1$：若 $\gcd(k,n)=d>1$，则 $a^k$ 的阶为 $n/d<n$，不能生成 $C_n$；若 $\gcd(k,n)=1$，由 Bézout 等式存在 $u,v$ 使 $uk+vn=1$，故 $a=(a^k)^u\cdot(a^n)^v=(a^k)^u\in\langle a^k\rangle$。

接着群结构。$\sigma_k\circ\sigma_l(a)=\sigma_k(a^l)=(a^k)^l=a^{kl}=\sigma_{kl}(a)$，故 $\sigma_k\circ\sigma_l=\sigma_{kl}$。映射 $k\mapsto\sigma_k$ 是 $(\ZZ/n\ZZ)^\times$ 到 $\operatorname{Aut}(C_n)$ 的同构。
\end{proof}

\medskip
\noindent\emph{小阶循环群的自同构。}
\begin{itemize}
\item $n=2$：$(\ZZ/2\ZZ)^\times=\{1\}$，$\operatorname{Aut}(C_2)$ 平凡（$C_2$ 只有唯一的非单位元，必须映到自身）。
\item $n=3$：$(\ZZ/3\ZZ)^\times=\{1,2\}\cong C_2$，$\operatorname{Aut}(C_3)\cong C_2$（唯一非平凡自同构 $a\mapsto a^2=a^{-1}$）。
\item $n=4$：$(\ZZ/4\ZZ)^\times=\{1,3\}\cong C_2$，$\operatorname{Aut}(C_4)\cong C_2$。
\item $n=5$：$(\ZZ/5\ZZ)^\times=\{1,2,3,4\}\cong C_4$（$2$ 是模 $5$ 的原根）。
\item $n=8$：$(\ZZ/8\ZZ)^\times=\{1,3,5,7\}\cong C_2\times C_2$（Klein 四元群），因 $3^2\equiv5^2\equiv7^2\equiv1\pmod{8}$。
\end{itemize}
\par\medskip

\begin{definition}[特征子群]
子群 $H\leqslant G$ 称为\textbf{特征子群}（记 $H\operatorname{char}G$），若对所有 $\varphi\in\operatorname{Aut}(G)$ 都有 $\varphi(H)=H$。
\end{definition}

特征子群比正规子群更强：正规子群只对内自同构不变，特征子群对所有自同构不变。

\begin{proposition}[特征子群的性质]
\label{prop:characteristic-subgroup}
\begin{enumerate}
\item 若 $H\operatorname{char}G$，则 $H\mathrel{\trianglelefteq} G$。
\item 若 $H\operatorname{char}K$ 且 $K\operatorname{char}G$，则 $H\operatorname{char}G$（传递性）。
\item 若 $H\operatorname{char}K$ 且 $K\mathrel{\trianglelefteq} G$，则 $H\mathrel{\trianglelefteq} G$。
\item $Z(G)\operatorname{char}G$，$G'\operatorname{char}G$（导群是特征子群）。
\end{enumerate}
\end{proposition}

\begin{proof}
(1) 内自同构是自同构的特例，故 $\varphi(H)=H$ 对所有 $\varphi\in\operatorname{Aut}(G)$ 成立蕴含 $c_g(H)=H$ 对所有 $g\in G$ 成立，即 $gHg^{-1}=H$，故 $H\mathrel{\trianglelefteq} G$。

(2) 对任意 $\varphi\in\operatorname{Aut}(G)$，$\varphi|_K\in\operatorname{Aut}(K)$（因 $K\operatorname{char}G$ 保证 $\varphi(K)=K$），故 $\varphi(H)=\varphi|_K(H)=H$（因 $H\operatorname{char}K$）。

(3) 对任意 $g\in G$，$c_g\in\operatorname{Inn}(G)\subseteq\operatorname{Aut}(G)$，$c_g(K)=K$（因 $K\mathrel{\trianglelefteq} G$），故 $c_g|_K\in\operatorname{Aut}(K)$，$c_g(H)=c_g|_K(H)=H$（因 $H\operatorname{char}K$）。

(4) 对任意 $\varphi\in\operatorname{Aut}(G)$ 和 $z\in Z(G)$：$\varphi(z)x\varphi(z)^{-1}=\varphi(z\varphi^{-1}(x)z^{-1})=\varphi(\varphi^{-1}(x))=x$，故 $\varphi(z)\in Z(G)$，即 $\varphi(Z(G))\subseteq Z(G)$；同理论 $\varphi^{-1}$ 得反向包含。导群 $G'=[G,G]$ 由换位子生成，$\varphi([x,y])=[\varphi(x),\varphi(y)]$，故 $\varphi(G')=G'$。
\end{proof}

\begin{theorem}[$N/C$ 定理]
\label{thm:nc-theorem}
设 $H\leqslant G$，则 $N_G(H)/C_G(H)$ 同构于 $\operatorname{Aut}(H)$ 的一个子群：
\[
N_G(H)/C_G(H)\hookrightarrow\operatorname{Aut}(H),
\]
其中 $N_G(H)=\{g\in G\mid gHg^{-1}=H\}$ 是正规化子，$C_G(H)=\{g\in G\mid gh=hg,\forall h\in H\}$ 是中心化子。特别地，$|N_G(H)/C_G(H)|\bigm||\operatorname{Aut}(H)|$。
\end{theorem}

\begin{proof}
定义映射 $\Phi:N_G(H)\to\operatorname{Aut}(H)$，$\Phi(g)=c_g|_H$（$h\mapsto ghg^{-1}$）。对 $g\in N_G(H)$，$c_g(H)=gHg^{-1}=H$，故 $c_g|_H$ 是 $H$ 到自身的同构，即 $\Phi(g)\in\operatorname{Aut}(H)$。

同态性。$\Phi(gh)(x)=ghxh^{-1}g^{-1}=g(hxh^{-1})g^{-1}=\Phi(g)\circ\Phi(h)(x)$。

核。$\Phi(g)=\id(H)$ 当且仅当 $ghg^{-1}=h$ 对所有 $h\in H$，即 $g\in C_G(H)$。故 $\operatorname{Ker}\Phi=C_G(H)$。

由第一同构定理，$N_G(H)/C_G(H)\cong\operatorname{Im}\Phi\leqslant\operatorname{Aut}(H)$。
\end{proof}

$N/C$ 定理在有限群论中极为有用：它把正规化子/中心化子的商群嵌入 $\operatorname{Aut}(H)$，给出了群结构的强约束。例如，若 $H$ 的自同构群很小（如 $H$ 是循环群且 $\varphi(|H|)$ 小），则 $N_G(H)/C_G(H)$ 也很小，这限制了 $H$ 在 $G$ 中的"非中心"部分。

\medskip
\noindent\emph{$N/C$ 定理应用于 Sylow 子群。}
设 $P$ 是 $G$ 的 Sylow $p$-子群。由 Sylow 定理，$N_G(P)$ 控制 $P$ 的共轭类。$N/C$ 定理给出 $N_G(P)/C_G(P)\hookrightarrow\operatorname{Aut}(P)$。若 $P\cong C_p$（$p$ 阶循环群），则 $|\operatorname{Aut}(P)|=p-1$，故 $|N_G(P)/C_G(P)|\bigm|(p-1)$。又 $P\leqslant C_G(P)$（$P$ 是 Abel 群时 $P$ 中心化自身），故 $|N_G(P)/P|\bigm|(p-1)$。这限制了 Sylow $p$-子群的正规化子的大小。

特别地，若 $p$ 是 $|G|$ 的最高次幂且 $p\nmid(p-1)$（总是如此），则 $N_G(P)/C_G(P)$ 的阶既整除 $p-1$ 又是 $p$ 的幂（因 $N_G(P)/P$ 的阶整除 $|G|/|P|$ 与 $p$ 互质），故 $N_G(P)=C_G(P)$，即 Sylow $p$-子群在自身正规化子中是中心的——这就是 Burnside 正规 $p$-补定理的一个关键推论。
\par\medskip

\begin{theorem}[Frattini 论证]
\label{thm:frattini-argument}
设 $G$ 是有限群，$N\mathrel{\trianglelefteq} G$，$P$ 是 $N$ 的 Sylow $p$-子群。则
\[
G=N\cdot N_G(P)=\{ng\mid n\in N,\,g\in N_G(P)\}.
\]
\end{theorem}

\begin{proof}
对任意 $g\in G$，$gPg^{-1}\leqslant gNg^{-1}=N$（因 $N\mathrel{\trianglelefteq} G$），且 $|gPg^{-1}|=|P|$，故 $gPg^{-1}$ 也是 $N$ 的 Sylow $p$-子群。由 Sylow 定理（在 $N$ 中应用），存在 $n\in N$ 使 $n(gPg^{-1})n^{-1}=P$，即 $(ng)P(ng)^{-1}=P$，故 $ng\in N_G(P)$，即 $g=n^{-1}\cdot ng\in N\cdot N_G(P)$。
\end{proof}

Frattini 论证看似简单，但它是有限群论中最常用的技巧之一。它的典型用法是：若 $N\mathrel{\trianglelefteq} G$ 且 $N$ 的 Sylow 子群 $P$ 的正规化子 $N_G(P)$ 已知（例如 $P$ 是 $N$ 的特征子群，则 $N_G(P)\geqslant N$），则可以分解 $G$ 并推导结构信息。

\medskip
\noindent\emph{Frattini 论证的应用：$p$-群的作用。}
设 $P$ 是 $G$ 的正规 Sylow $p$-子群（即 $P\mathrel{\trianglelefteq} G$）。在 Frattini 论证中取 $N=P$，$P$ 自身是 $P$ 的唯一 Sylow $p$-子群，故 $N_G(P)=G$。论证给出 $G=P\cdot G$（平凡），但更有趣的是取 $N=P$ 的某个正规子群 $K\mathrel{\trianglelefteq} G$，$K\leqslant P$：$G=K\cdot N_G(Q)$，其中 $Q$ 是 $K$ 的 Sylow $p$-子群。若 $K$ 是 $p$-群，$Q=K$，$N_G(Q)\geqslant N_G(K)=G$，故 $G=K\cdot G$（仍平凡）。真正的力量体现在 $K$ 不是 $p$-群时：Frattini 论证把 $G$ 分解为 $K$ 和一个可能更小的群 $N_G(Q)$，可以用于归纳证明。
\par\medskip

\begin{proposition}[中心化子的正规化子性质]
\label{prop:centralizer-normalizer}
设 $H\leqslant G$。则：
\begin{enumerate}
\item $C_G(H)\mathrel{\trianglelefteq} N_G(H)$。
\item 若 $H$ 是 Abel 群，则 $H\leqslant C_G(H)$，且 $H\mathrel{\trianglelefteq} N_G(H)$。
\item 若 $H\mathrel{\trianglelefteq} G$ 且 $H$ 是 Abel 群，则 $H\leqslant Z(G)$ 当且仅当 $C_G(H)=G$。
\end{enumerate}
\end{proposition}

\begin{proof}
(1) 对 $n\in N_G(H)$ 和 $c\in C_G(H)$，需证 $ncn^{-1}\in C_G(H)$。对任意 $h\in H$：
\[
ncn^{-1}h(ncn^{-1})^{-1}=nc(n^{-1}hn)c^{-1}n^{-1}.
\]
因 $n\in N_G(H)$，$n^{-1}hn\in H$，而 $c\in C_G(H)$ 与 $H$ 的所有元素交换，故 $c(n^{-1}hn)c^{-1}=n^{-1}hn$。于是上式 $=n(n^{-1}hn)n^{-1}=h$。

(2) $H$ Abel 蕴含 $H$ 的元素彼此交换，故 $H\leqslant C_G(H)$。由 (1) $C_G(H)\mathrel{\trianglelefteq} N_G(H)$，故 $H\leqslant C_G(H)\mathrel{\trianglelefteq} N_G(H)$，特别 $H\mathrel{\trianglelefteq} N_G(H)$（子群的正规子群不一定正规于更大的群，但这里 $H\leqslant C_G(H)\mathrel{\trianglelefteq} N_G(H)$ 且 $H\mathrel{\trianglelefteq} C_G(H)$（因 $C_G(H)$ 中 $H$ 是中心的），由正规性的传递性不成立的一般警告，需直接验证：对 $n\in N_G(H)$，$nHn^{-1}=H$，这正是 $N_G(H)$ 的定义）。

(3) $H\leqslant Z(G)$ 当且仅当 $G$ 的每个元素与 $H$ 的每个元素交换，即 $G\leqslant C_G(H)$，即 $C_G(H)=G$。
\end{proof}

\begin{theorem}[对称群的自同构群]
\label{thm:symmetric-aut}
对 $n\geqslant7$，$n\neq6$，
\[
\operatorname{Aut}(S_n)=\operatorname{Inn}(S_n)\cong S_n,
\]
即 $S_n$ 的所有自同构都是内自同构（$\operatorname{Out}(S_n)=1$）。唯一的例外是 $n=6$：$\operatorname{Out}(S_6)\cong C_2$，$S_6$ 有外自同构。
\end{theorem}

这个定理的证明较长（关键是证明每个自同构把对换映到对换，然后利用对换生成 $S_n$），但结论极为重要：$S_n$（$n\neq6$）的对称性"全部来自自身"——没有外部对称。$n=6$ 的外自同构是有限群论中最著名的例外之一，它与 $S_6$ 的特殊结构有关（$S_6$ 有 12 个 Sylow $5$-子群而非"通常"的 6 个，导致额外的对称性）。

\medskip
\noindent\emph{$S_6$ 的外自同构。}
$S_6$ 的外自同构 $\sigma\in\operatorname{Aut}(S_6)\setminus\operatorname{Inn}(S_6)$ 把对换（2-循环）映到\emph{三对对换的乘积}（即 $(ab)\mapsto(a'b')(c'd')(e'f')$，将 6 个字母分成三对）。这种映射不来自 $S_6$ 的任何共轭，但保持了 $S_6$ 的乘法结构。$S_6$ 的外自同构与数学中许多"偶然例外"现象相关，如 Cremona 变换、正二十面体的特殊对称等。
\par\medskip
